% Options for packages loaded elsewhere
\PassOptionsToPackage{unicode}{hyperref}
\PassOptionsToPackage{hyphens}{url}
\PassOptionsToPackage{dvipsnames,svgnames,x11names}{xcolor}
%
\documentclass[
  11pt,
]{article}

\usepackage{amsmath,amssymb}
\usepackage{iftex}
\ifPDFTeX
  \usepackage[T1]{fontenc}
  \usepackage[utf8]{inputenc}
  \usepackage{textcomp} % provide euro and other symbols
\else % if luatex or xetex
  \usepackage{unicode-math}
  \defaultfontfeatures{Scale=MatchLowercase}
  \defaultfontfeatures[\rmfamily]{Ligatures=TeX,Scale=1}
\fi
\usepackage{lmodern}
\ifPDFTeX\else  
    % xetex/luatex font selection
\fi
% Use upquote if available, for straight quotes in verbatim environments
\IfFileExists{upquote.sty}{\usepackage{upquote}}{}
\IfFileExists{microtype.sty}{% use microtype if available
  \usepackage[]{microtype}
  \UseMicrotypeSet[protrusion]{basicmath} % disable protrusion for tt fonts
}{}
\makeatletter
\@ifundefined{KOMAClassName}{% if non-KOMA class
  \IfFileExists{parskip.sty}{%
    \usepackage{parskip}
  }{% else
    \setlength{\parindent}{0pt}
    \setlength{\parskip}{6pt plus 2pt minus 1pt}}
}{% if KOMA class
  \KOMAoptions{parskip=half}}
\makeatother
\usepackage{xcolor}
\usepackage[margin=1in]{geometry}
\setlength{\emergencystretch}{3em} % prevent overfull lines
\setcounter{secnumdepth}{-\maxdimen} % remove section numbering
% Make \paragraph and \subparagraph free-standing
\ifx\paragraph\undefined\else
  \let\oldparagraph\paragraph
  \renewcommand{\paragraph}[1]{\oldparagraph{#1}\mbox{}}
\fi
\ifx\subparagraph\undefined\else
  \let\oldsubparagraph\subparagraph
  \renewcommand{\subparagraph}[1]{\oldsubparagraph{#1}\mbox{}}
\fi


\providecommand{\tightlist}{%
  \setlength{\itemsep}{0pt}\setlength{\parskip}{0pt}}\usepackage{longtable,booktabs,array}
\usepackage{calc} % for calculating minipage widths
% Correct order of tables after \paragraph or \subparagraph
\usepackage{etoolbox}
\makeatletter
\patchcmd\longtable{\par}{\if@noskipsec\mbox{}\fi\par}{}{}
\makeatother
% Allow footnotes in longtable head/foot
\IfFileExists{footnotehyper.sty}{\usepackage{footnotehyper}}{\usepackage{footnote}}
\makesavenoteenv{longtable}
\usepackage{graphicx}
\makeatletter
\def\maxwidth{\ifdim\Gin@nat@width>\linewidth\linewidth\else\Gin@nat@width\fi}
\def\maxheight{\ifdim\Gin@nat@height>\textheight\textheight\else\Gin@nat@height\fi}
\makeatother
% Scale images if necessary, so that they will not overflow the page
% margins by default, and it is still possible to overwrite the defaults
% using explicit options in \includegraphics[width, height, ...]{}
\setkeys{Gin}{width=\maxwidth,height=\maxheight,keepaspectratio}
% Set default figure placement to htbp
\makeatletter
\def\fps@figure{htbp}
\makeatother

\usepackage{array}
\usepackage{booktabs}
\usepackage{graphicx}
\usepackage{longtable}
\usepackage{needspace}
\usepackage{caption}
\usepackage{float}
\usepackage{fancyhdr}
\pagestyle{fancy}
\fancyhf{}
\fancyhead[L]{Video Quality Analysis Report}
\fancyhead[R]{Fifth Avenue Studios}
\fancyfoot[L]{\today}
\fancyfoot[C]{\thepage}
\fancyfoot[R]{Gil Raitses}
\setlength{\extrarowheight}{1pt}
\renewcommand{\arraystretch}{1.0}
\makeatletter
\makeatother
\makeatletter
\makeatother
\makeatletter
\@ifpackageloaded{caption}{}{\usepackage{caption}}
\AtBeginDocument{%
\ifdefined\contentsname
  \renewcommand*\contentsname{Table of contents}
\else
  \newcommand\contentsname{Table of contents}
\fi
\ifdefined\listfigurename
  \renewcommand*\listfigurename{List of Figures}
\else
  \newcommand\listfigurename{List of Figures}
\fi
\ifdefined\listtablename
  \renewcommand*\listtablename{List of Tables}
\else
  \newcommand\listtablename{List of Tables}
\fi
\ifdefined\figurename
  \renewcommand*\figurename{Figure}
\else
  \newcommand\figurename{Figure}
\fi
\ifdefined\tablename
  \renewcommand*\tablename{Table}
\else
  \newcommand\tablename{Table}
\fi
}
\@ifpackageloaded{float}{}{\usepackage{float}}
\floatstyle{ruled}
\@ifundefined{c@chapter}{\newfloat{codelisting}{h}{lop}}{\newfloat{codelisting}{h}{lop}[chapter]}
\floatname{codelisting}{Listing}
\newcommand*\listoflistings{\listof{codelisting}{List of Listings}}
\makeatother
\makeatletter
\@ifpackageloaded{caption}{}{\usepackage{caption}}
\@ifpackageloaded{subcaption}{}{\usepackage{subcaption}}
\makeatother
\makeatletter
\@ifpackageloaded{tcolorbox}{}{\usepackage[skins,breakable]{tcolorbox}}
\makeatother
\makeatletter
\@ifundefined{shadecolor}{\definecolor{shadecolor}{rgb}{.97, .97, .97}}
\makeatother
\makeatletter
\makeatother
\makeatletter
\makeatother
\ifLuaTeX
  \usepackage{selnolig}  % disable illegal ligatures
\fi
\IfFileExists{bookmark.sty}{\usepackage{bookmark}}{\usepackage{hyperref}}
\IfFileExists{xurl.sty}{\usepackage{xurl}}{} % add URL line breaks if available
\urlstyle{same} % disable monospaced font for URLs
\hypersetup{
  pdftitle={Video Quality Analysis Report},
  colorlinks=true,
  linkcolor={blue},
  filecolor={Maroon},
  citecolor={Blue},
  urlcolor={Blue},
  pdfcreator={LaTeX via pandoc}}

\title{Video Quality Analysis Report}
\usepackage{etoolbox}
\makeatletter
\providecommand{\subtitle}[1]{% add subtitle to \maketitle
  \apptocmd{\@title}{\par {\large #1 \par}}{}{}
}
\makeatother
\subtitle{Flaco Owl Documentary Footage - Fifth Avenue Studios}
\author{}
\date{2025-11-13}

\begin{document}
\maketitle
\begin{abstract}
This report analyzes 14.06 seconds of video footage documenting Flaco
the owl's first landing on Fifth Avenue on February 4, 2024, the night
he was freed from the Central Park Zoo. This clip distinguishes itself
with the highest valuation rate (\$210.94/second) and highest estimated
value range (\$2,967 - \$3,223) among the three clips, driven by
superior technical quality metrics including the best exposure
(98.0/100), highest smoothness (0.955), and minimal color compression
(0.938) of the set. Technical analysis reveals Full HD resolution (1080
× 1920 pixels, 30 fps) with overall quality score of 51.8/100, requiring
post-processing for sharpness (33.15) and stability (4.55). Unlike the
other clips, this footage shows no subject detection in analyzed frames,
making it suitable for establishing shots or atmospheric context. The
footage achieves 100\% usability (14.06 seconds usable) with 17.2\%
meeting high-quality standards (2.43 seconds). Narrative significance
scores 90/100 due to the rare/unique event nature, historic moment value
and multiple-use potential throughout documentary production. This clip
offers the highest per-second value and best technical foundation for
post-processing enhancement.
\end{abstract}
\ifdefined\Shaded\renewenvironment{Shaded}{\begin{tcolorbox}[borderline west={3pt}{0pt}{shadecolor}, frame hidden, interior hidden, boxrule=0pt, sharp corners, breakable, enhanced]}{\end{tcolorbox}}\fi

\needspace{8cm}
\section{Executive Summary}

This report provides a technical and narrative analysis of video footage
documenting Flaco the owl's first-night landing on Fifth Avenue after
being freed from the Central Park Zoo. The footage has been evaluated
for documentary production suitability and content significance.

\begin{table}[h]
\centering
\small
\caption{Executive Summary Findings}
\label{tab-exec-summary}
\begin{tabular}{p{3cm}p{8cm}}
\toprule
\textbf{Finding} & \textbf{Value} \\
\midrule
Total Duration & 14.06 seconds \\
Usable Footage & 14.06 seconds (100\%) \\
High Quality Segments & 2.43 seconds (17.2\%) \\
Narrative Significance & 90/100 (high importance) \\
Estimated Value Range & \$2,967 - \$3,223 \\
\bottomrule
\end{tabular}
\end{table}

\begin{figure}[h]
\centering
\includegraphics[width=0.85\textwidth]{figures_3787/value_range.pdf}
\caption{This clip has the highest valuation range (\$2,967 - \$3,223) and highest per-second rate (\$210.94) among the three clips, reflecting superior technical quality multipliers. Bar chart comparing Conservative (\$2,967), Recommended and Optimistic (\$3,223) valuation estimates.}
\label{fig-value-range}
\end{figure}

\needspace{8cm}
\section{Technical Specifications}

\begin{table}[h]
\centering
\small
\caption{Technical Specifications}
\label{tab-tech-specs}
\begin{tabular}{p{3cm}p{7.5cm}}
\toprule
\textbf{Property} & \textbf{Value} \\
\midrule
Resolution & 1080 × 1920 pixels (Full HD 1080p, Portrait) \\
Frame Rate & 30.00 fps \\
Duration & 14.06 seconds \\
Total Frames & 422 frames \\
File Size & 23.09 MB \\
File Format & QuickTime (.MOV) \\
Video Codec & H.264 (AVC) \\
Color Depth & 8-bit \\
Pixel Format & YUV 4:2:0 \\
Aspect Ratio & 9:16 (Portrait) \\
\bottomrule
\end{tabular}
\end{table}

\needspace{10cm}
\section{Visual Quality Metrics}

\subsection{Core Quality Metrics}

\begin{table}[h]
\centering
\small
\caption{Core Quality Metrics}
\label{tab-quality-metrics}
\begin{tabular}{p{3.5cm}p{2cm}p{5.5cm}}
\toprule
\textbf{Metric} & \textbf{Score} & \textbf{Assessment} \\
\midrule
Sharpness & 33.15 & Very Soft/Blurry \\
Brightness & 72.0/255 & Underexposed \\
Contrast & 43.89 & High Contrast \\
Noise Level & 5.68 & Very Clean \\
Exposure & 98.0/100 & Good Exposure \\
Stability & 4.55 & Unstable/Shaky \\
Smoothness & 0.955 & Very Smooth \\
Color Compression & 0.938 & Minimal Compression \\
\bottomrule
\end{tabular}
\end{table}

\needspace{12cm}
\subsection{Quality Assessment}

\textbf{Strengths:}

\begin{table}[h]
\centering
\small
\caption{Quality Strengths}
\label{tab-strengths}
\begin{tabular}{p{2cm}p{7cm}}
\toprule
\textbf{Strength} & \textbf{Description} \\
\midrule
Resolution & High resolution (Full HD 1080p) suitable for broadcast \\
Exposure & Well-exposed footage with minimal over/under exposure \\
Noise Level & Clean footage with minimal noise (5.68) \\
Smoothness & High temporal smoothness (0.955) \\
Compression & Minimal color compression artifacts (0.938) \\
\bottomrule
\end{tabular}
\end{table}

\textbf{Post-Processing Considerations:}

\begin{table}[h]
\centering
\small
\caption{Post-Processing Requirements}
\label{tab-post-processing}
\begin{tabular}{p{2cm}p{4cm}p{3cm}}
\toprule
\textbf{Issue} & \textbf{Description} & \textbf{Required Action} \\
\midrule
Sharpness & Footage is soft/blurry (sharpness: 33.15) & Requires sharpening \\
Stability & Camera shake present (stability: 4.55) & Requires stabilization \\
\bottomrule
\end{tabular}
\end{table}

\begin{figure}[h]
\centering
\includegraphics[width=0.7\textwidth]{figures_3787/quality_metrics_radar.pdf}
\caption{This clip shows the best exposure and smoothness scores among the three clips, making it ideal for post-processing enhancement. Radar chart showing normalized quality scores across six dimensions with exposure (98.0) and smoothness (95.5) as standout metrics.}
\label{fig-quality-radar}
\end{figure}

\needspace{22cm}
\section{Content Significance Analysis}

\subsection{Subject Detection}

\begin{table}[h]
\centering
\small
\caption{Subject Detection Metrics}
\label{tab-subject-detection}
\begin{tabular}{p{4.5cm}p{5.5cm}}
\toprule
\textbf{Metric} & \textbf{Value} \\
\midrule
Frames Analyzed & 29 frames \\
Frames with Subject (Bird/Owl) & 0 frames (0.0\%) \\
Frames with Dog & 0 frames \\
Frames with Interaction & 0 frames \\
Subject Visible Throughout & No \\
Interaction Moments Detected & No \\
Rare/Unique Event & Yes \\
\bottomrule
\end{tabular}
\end{table}

\subsection{Event Significance}

This footage documents a rare and historic event: Flaco the owl's first
landing on Fifth Avenue on the night he was freed from the Central Park
Zoo. This represents an important moment in the narrative arc of the
documentary, capturing the first public appearance of Flaco after his
escape.

\needspace{20cm}
\section{Narrative Significance (Story Value)}

\needspace{12cm}
\subsection{Narrative Score: 90/100}

\begin{table}[h]
\centering
\small
\caption{Narrative Significance Factors}
\label{tab-narrative-factors}
\begin{tabular}{p{3cm}p{6cm}}
\toprule
\textbf{Factor} & \textbf{Contribution} \\
\midrule
Rare/Unique Event & First-night landing on Fifth Ave \\
Historic Moment & First appearance after zoo escape \\
Subject Visibility & Subject not detected in analyzed frames \\
Multiple-Use Potential & Can be used multiple times throughout documentary \\
Story Importance & High importance \\
\bottomrule
\end{tabular}
\end{table}

\textbf{Multiple-Use Potential:} This footage can be strategically
reused throughout the documentary in:

\begin{table}[h]
\centering
\small
\caption{Multiple-Use Applications}
\label{tab-multiple-use}
\begin{tabular}{p{3.5cm}p{7.5cm}}
\toprule
\textbf{Use Case} & \textbf{Description} \\
\midrule
Opening Sequence & Can be used in documentary opening \\
Transitional Moments & Suitable for scene transitions \\
Key Narrative Scenes & Can feature in important story moments \\
Closing Sequences & Appropriate for documentary conclusion \\
\bottomrule
\end{tabular}
\end{table}

\begin{figure}[h]
\centering
\includegraphics[width=0.7\textwidth]{figures_3787/narrative_significance.pdf}
\caption{Narrative score of 90/100 reflects historic moment value despite no subject detection, making this clip valuable for establishing context and atmosphere. Pie chart showing narrative significance score of 90/100.}
\label{fig-narrative}
\end{figure}

\needspace{25cm}
\section{Usability Analysis}

\subsection{Usability Breakdown}

\begin{figure}[h]
\centering
\includegraphics[width=0.7\textwidth]{figures_3787/usability_breakdown.pdf}
\caption{Shortest clip (14.06s) with 100\% usability, making it a compact, production-ready segment. Bar chart comparing total duration (14.06s, 100\%), usable seconds (14.06s, 100\%) and high-quality seconds (2.43s, 17.2\%).}
\label{fig-usability-breakdown}
\end{figure}

\subsection{Per-Second Breakdown}

\begin{table}[h]
\centering
\small
\caption{Per-Second Breakdown}
\label{tab-per-second}
\begin{tabular}{p{2.5cm}cc}
\toprule
\textbf{Category} & \textbf{Seconds} & \textbf{Percentage} \\
\midrule
Total Duration & 14.06 & 100.0\% \\
Usable Seconds & 14.06 & 100.0\% \\
High Quality Seconds & 2.43 & 17.2\% \\
\bottomrule
\end{tabular}
\end{table}

\textbf{Usable Segments:} 1 continuous segment (0.00s - 14.06s)

\subsection{Clear, Focused Segments}

No continuous clear/focused segments of 1 second or longer were
identified. The footage requires post-processing (sharpening and
stabilization) to achieve optimal clarity for broadcast use.

\needspace{25cm}
\section{Valuation Analysis}

\subsection{Valuation Breakdown}

\begin{figure}[h]
\centering
\includegraphics[width=0.75\textwidth]{figures_3787/valuation_breakdown.pdf}
\caption{Highest per-second rate (\$210.94) among the three clips, driven by the best quality multiplier (0.75×) reflecting superior exposure, smoothness and compression metrics. Waterfall chart showing how base rate (\$50.00) is modified by sequential multipliers: Quality (0.75×), Resolution (1.00×), Content Significance (1.50×), Narrative Significance (2.50×) and Multiple-Use (1.50×).}
\label{fig-valuation-breakdown}
\end{figure}

\needspace{8cm}
\subsection{Rate Calculation}

\enlargethispage{2cm}
\vspace{-0.5cm}
\begin{table}[H]
\centering
\footnotesize
\caption{Rate Calculation}
\label{tab-rate-calculation}
\begin{tabular}{p{3.5cm}p{1.8cm}p{4.2cm}}
\toprule
\textbf{Factor} & \textbf{Multiplier} & \textbf{Value} \\
\midrule
Base Rate (per usable second) & 1.00× & \$50.00 \\
Quality Multiplier & 0.75× & \\
Resolution Multiplier & 1.00× & \\
Content Significance Multiplier & 1.50× & \\
Narrative Significance Multiplier & 2.50× & \\
Multiple-Use Multiplier & 1.50× & \\
\textbf{Final Rate per Second} & & \textbf{\$210.94} \\
\bottomrule
\end{tabular}
\end{table}

\needspace{8cm}
\subsection{Estimated Value Range}

\begin{table}[h]
\centering
\small
\caption{Estimated Value Range}
\label{tab-value-range}
\begin{tabular}{p{3cm}p{5cm}p{3cm}}
\toprule
\textbf{Estimate Type} & \textbf{Calculation} & \textbf{Value} \\
\midrule
Conservative & All usable seconds × \$210.94 & \$2,967 \\
Optimistic & Premium segments at 1.5× rate & \$3,223 \\
\textbf{Recommended Range} & & \textbf{\$2,967 - \$3,223} \\
\bottomrule
\end{tabular}
\end{table}

\needspace{12cm}
\subsection{Value Factors}

\begin{table}[h]
\centering
\small
\caption{Value Factors}
\label{tab-value-factors}
\begin{tabular}{p{3cm}p{7cm}}
\toprule
\textbf{Factor} & \textbf{Description} \\
\midrule
Usability & High usability (100\% usable) - most of the clip is production-ready \\
Archival Value & Rare/unique event adds archival and narrative value \\
Story Value & Story moment with reuse potential increases value \\
Multiple-Use & Footage can be used multiple times throughout documentary structure \\
\bottomrule
\end{tabular}
\end{table}

\section{Conclusion}

This footage represents an important moment in Flaco the owl's story,
documenting his first landing on Fifth Avenue after being freed from the
Central Park Zoo. This clip stands out among the three with the highest
valuation (\$2,967 - \$3,223) and highest per-second rate (\$210.94),
driven by superior technical metrics including the best exposure (98.0),
highest smoothness (0.955), and minimal color compression (0.938) of the
set. While sharpness and stability require post-processing, the strong
technical foundation makes this clip ideal for enhancement. Unlike the
other clips, no subject was detected in analyzed frames, positioning
this footage as valuable for establishing shots, atmospheric context, or
background elements. The narrative significance, rarity and multiple-use
potential, combined with the highest technical quality scores, make this
footage the most valuable of the three clips for documentary production.



\end{document}
